\documentclass[a4paper,13pt,3p,twoside]{report}
\usepackage{scrextend}
\changefontsizes{13pt}
\usepackage[utf8]{vietnam}
\usepackage[top=2cm, bottom=2cm, left=3.5cm, right=2.5cm]{geometry}
\usepackage{xurl}
\usepackage{appendix}
\usepackage{xcolor}
\usepackage{outlines}
\usepackage{graphicx} % Cho phép chèn hỉnh ảnh
\usepackage{fancybox} % Tạo khung box
\usepackage{indentfirst} % Thụt đầu dòng ở dòng đầu tiên trong đoạn
\usepackage{amsthm} % Cho phép thêm các môi trường định nghĩa
\usepackage{latexsym} % Các kí hiệu toán học
\usepackage{amsmath} % Hỗ trợ một số biểu thức toán học
\usepackage{amssymb} % Bổ sung thêm kí hiệu về toán học
\usepackage{amsbsy} % Hỗ trợ các kí hiệu in đậm
\usepackage{times} % Chọn font Time New Romans
\usepackage{array} % Tạo bảng array
\usepackage{enumitem} % Cho phép thay đổi kí hiệu của list
\usepackage{subfiles} % Chèn các file nhỏ, giúp chia các chapter ra nhiều file hơn
\usepackage{titlesec} % Giúp chỉnh sửa các tiêu đề, đề mục như chương, phần,..
\usepackage{titletoc}
\usepackage{chngcntr} % Dùng để thiết lập lại cách đánh số caption,..
\usepackage{pdflscape} % Đưa các bảng có kích thước đặt theo chiều ngang giấy
\usepackage{afterpage}
\usepackage[ruled,vlined]{algorithm2e}  % Hỗ trợ viết các giải thuật
\usepackage{capt-of} % Cho phép sử dụng caption lớn đối với landscape page
\usepackage{multirow} % Merge cells
\usepackage{longtable} % Bảng dài, ngắt qua trang (bổ sung so với mẫu)
\usepackage{fancyhdr} % Cho phép tùy biến header và footer

\usepackage[font=small,labelfont=bf]{caption}

\usepackage{listings}
\usepackage{float}
\usepackage[section]{placeins} % Bảng, hình không trôi qua ranh giới mục (bổ sung so với mẫu)
\raggedbottom % Không giãn khoảng cách đoạn khi bảng [H] nhảy trang (twoside mặc định \flushbottom)
\usepackage{subcaption}
\usepackage{tikz} % Sơ đồ quy trình và mặt bằng nút giao (bổ sung so với mẫu)
\usetikzlibrary{arrows.meta,positioning,calc,shapes.geometric,fit,backgrounds,patterns}

\usepackage[nonumberlist, nopostdot, nogroupskip, acronym]{glossaries}
\usepackage{glossary-superragged}
\setglossarystyle{superraggedheaderborder}
\usepackage{setspace}
\usepackage{parskip}

% package content table
\usepackage{tocbasic}

% ===================================================

% Mẫu dùng backend=bibtex; với style=ieee, bibtex luôn báo "Please (re)run BibTeX", nên dùng biber
\usepackage[backend=biber,style=ieee]{biblatex}

\addbibresource{\subfix{reference.bib}} % \subfix: đúng đường dẫn cả khi biên dịch riêng một chương
% URL dài trong danh mục tài liệu được ngắt ở mọi ký tự, tránh tràn lề
\setcounter{biburlnumpenalty}{9000}
\setcounter{biburlucpenalty}{9000}
\setcounter{biburllcpenalty}{9000}

% Tránh tràn lề: cho phép giãn thêm khi dàn dòng, và cho ngắt dòng sau dấu gạch dưới
% trong tên định danh như \texttt{queue\_space\_seeking}
\setlength{\emergencystretch}{3em}
\let\gachduoi\_
\renewcommand{\_}{\gachduoi\allowbreak}

%\makeglossaries
\makenoidxglossaries

% Danh mục thuật ngữ và từ viết tắt
\newglossaryentry{GEH}{
    type=\acronymtype,
    name={GEH},
    description={Chỉ số so sánh lưu lượng mô phỏng và quan sát (Geoffrey E. Havers statistic)},
    first={chỉ số GEH}
}
\newglossaryentry{LLM}{
    type=\acronymtype,
    name={LLM},
    description={Mô hình ngôn ngữ lớn (Large Language Model)},
    first={mô hình ngôn ngữ lớn (Large Language Model - LLM)}
}
\newglossaryentry{OSM}{
    type=\acronymtype,
    name={OSM},
    description={Dữ liệu bản đồ mở OpenStreetMap},
    first={OpenStreetMap (OSM)}
}
\newglossaryentry{PET}{
    type=\acronymtype,
    name={PET},
    description={Thời gian sau xâm phạm (Post-Encroachment Time)},
    first={thời gian sau xâm phạm (Post-Encroachment Time - PET)}
}
\newglossaryentry{SUMO}{
    type=\acronymtype,
    name={SUMO},
    description={Trình mô phỏng giao thông vi mô mã nguồn mở (Simulation of Urban MObility)},
    first={Simulation of Urban MObility (SUMO)}
}
\newglossaryentry{TraCI}{
    type=\acronymtype,
    name={TraCI},
    description={Giao diện điều khiển mô phỏng SUMO khi đang chạy (Traffic Control Interface)},
    first={Traffic Control Interface (TraCI)}
}
% Bổ sung theo phản biện 22/11/2026 (HB-15, ND-11, VP-16): viết tắt dùng trong thân bài mà danh mục chưa có
\newglossaryentry{GPU}{
    type=\acronymtype,
    name={GPU},
    description={Bộ xử lý đồ họa (Graphics Processing Unit)}
}
\newglossaryentry{IDM}{
    type=\acronymtype,
    name={IDM},
    description={Mô hình bám xe người lái thông minh (Intelligent Driver Model)}
}
\newglossaryentry{IQR}{
    type=\acronymtype,
    name={IQR},
    description={Khoảng tứ phân vị (Interquartile Range)}
}
\newglossaryentry{JSON}{
    type=\acronymtype,
    name={JSON},
    description={Định dạng văn bản cho dữ liệu có cấu trúc (JavaScript Object Notation)}
}
\newglossaryentry{MCP}{
    type=\acronymtype,
    name={MCP},
    description={Giao thức kết nối LLM với công cụ ngoài (Model Context Protocol)}
}
\newglossaryentry{MOBIL}{
    type=\acronymtype,
    name={MOBIL},
    description={Mô hình chuyển làn MOBIL (Minimizing Overall Braking Induced by Lane changes)}
}
\newglossaryentry{SSM}{
    type=\acronymtype,
    name={SSM},
    description={Chỉ số an toàn thay thế (Surrogate Safety Measures), thiết bị đo xung đột của SUMO}
}
\newglossaryentry{TCCS}{
    type=\acronymtype,
    name={TCCS},
    description={Tiêu chuẩn cơ sở}
}
\newglossaryentry{TCVN}{
    type=\acronymtype,
    name={TCVN},
    description={Tiêu chuẩn quốc gia Việt Nam}
}
\newglossaryentry{TLS}{
    type=\acronymtype,
    name={TLS},
    description={Hệ đèn tín hiệu giao thông trong SUMO (Traffic Light System)}
}
\newglossaryentry{TPE}{
    type=\acronymtype,
    name={TPE},
    description={Bộ lấy mẫu của Optuna dựa trên ước lượng mật độ (Tree-structured Parzen Estimator)}
}
\newglossaryentry{vType}{
    type=\acronymtype,
    name={vType},
    description={Kiểu phương tiện trong SUMO, tập thuộc tính hành vi của một loại xe (vehicle type)}
}
\newglossaryentry{XSD}{
    type=\acronymtype,
    name={XSD},
    description={Lược đồ định nghĩa cấu trúc tệp XML (XML Schema Definition)}
}


% Giá trị tạm thời (cấu hình thay thế, mục tiêu thay thế, tham số giả định): in màu đỏ kèm dấu dagger
% để vẫn nhận ra khi in đen trắng; sẽ thay bằng số đo thực
\newcommand{\tamthoi}[1]{\textcolor{red}{#1}\ensuremath{^{\dagger}}}

% Ghi chú làm việc (TODO, câu hỏi cho GVHD): in màu xanh trong bản nháp.
% Bản nộp: \bannoptrue ẩn toàn bộ ghi chú.
\newif\ifbannop
\bannoptrue
\newcommand{\ghichu}[1]{\ifbannop\else\textcolor{blue}{[#1]}\fi}

% ===================================================


\fancypagestyle{plain}{%
\fancyhf{} % clear all header and footer fields
\fancyfoot[RO,RE]{\thepage} %RO=right odd, RE=right even
\renewcommand{\headrulewidth}{0pt}
\renewcommand{\footrulewidth}{0pt}}

\setlength{\headheight}{18pt} % fancyhdr yêu cầu tối thiểu ~17.98pt với cỡ chữ 13pt

\def \TITLE{LUẬN VĂN THẠC SĨ}
\def \AUTHOR{QUYỀN ĐỨC CÔNG}

% ===================================================
\titleformat
    {\chapter} % command
    [hang] % shape
    {\centering\bfseries} % format
    {CHƯƠNG \thechapter.\ } % label
    {0pt} %sep
    {} % before
    [] % after
\titlespacing*{\chapter}{0pt}{-20pt}{20pt}

\titleformat
    {\section} % command
    [hang] % shape
    {\bfseries} % format
    {\thechapter.\arabic{section}\ \ \ \ } % label
    {0pt} %sep
    {} % before
    [] % after
\titlespacing{\section}{0pt}{\parskip}{0.5\parskip}

\titleformat
    {\subsection} % command
    [hang] % shape
    {\bfseries} % format
    {\thechapter.\arabic{section}.\arabic{subsection}\ \ \ \ } % label
    {0pt} %sep
    {} % before
    [] % after
\titlespacing{\subsection}{30pt}{\parskip}{0.5\parskip}

\renewcommand\thesubsubsection{\alph{subsubsection}}
\titleformat
    {\subsubsection} % command
    [hang] % shape
    {\bfseries} % format
    {\alph{subsubsection}, \ } % label
    {0pt} %sep
    {} % before
    [] % after
\titlespacing{\subsubsection}{50pt}{\parskip}{0.5\parskip}

% ===================================================
\usepackage{hyperref}
\hypersetup{pdfborder = {0 0 0}}
\hypersetup{pdftitle={\TITLE},
	pdfauthor={\AUTHOR}}

\usepackage[all]{hypcap} % Cho phép tham chiếu chính xác đến hình ảnh và bảng biểu

\graphicspath{{\subfix{./}}} % Đường dẫn hình viết dạng Figure/<tệp> tính từ thư mục gốc

\counterwithin{figure}{chapter} % Đánh số hình ảnh kèm theo chapter. Ví dụ: Hình 1.1, 1.2,..

\title{\bf \TITLE}
\author{\AUTHOR}

\setcounter{secnumdepth}{3} % Cho phép subsubsection trong report

\theoremstyle{definition}
\newtheorem{example}{Ví dụ}[chapter] % Định nghĩa môi trường ví dụ

\onehalfspacing
\setlength{\parskip}{6pt}
\setlength{\parindent}{15pt}

% Nhãn tiếng Việt (mẫu SOICT bản tiếng Anh; nội dung luận văn viết tiếng Việt)
\renewcommand{\figurename}{Hình}
\renewcommand{\tablename}{Bảng}
\renewcommand{\chaptername}{CHƯƠNG}


% =========================== BODY ===============
\begin{document}
\subfile{Cover} % Phần bìa
\subfile{Cover2} % Phần bìa lót

% ===================================================
\pagestyle{empty} % Header và footer rỗng

\newpage
\pagenumbering{gobble}
\subfile{Chapter/0_2_Acknowledgment.tex}

\newpage
\pagenumbering{gobble} % Xóa page numbering ở cuối trang
\subfile{Chapter/0_3_Abstract.tex}

\newpage
\pagenumbering{gobble}
\subfile{Chapter/0_4_Abstract_English.tex}


% ===================================================
\newpage
\pagenumbering{gobble} % Xóa page numbering ở cuối trang
\renewcommand*\contentsname{MỤC LỤC}
\contentsmargin{2.5em} % Chừa chỗ cho số trang ba chữ số
\titlecontents{chapter}
    [0.0cm]             % left margin
    {\bfseries\vspace{0.3cm}}                  % above code
    {{\bfseries{\scshape} CHƯƠNG \thecontentslabel.\ }} % numbered format
    {}         % unnumbered format
    {\titlerule*[0.3pc]{.}\contentspage}         % filler-page-format, e.g dots

\titlecontents{section}
    [0.0cm]             % left margin
    {\vspace{0.3cm}}                  % above code
    {\thecontentslabel \ } % numbered format
    {}         % unnumbered format
    {\titlerule*[0.3pc]{.}\contentspage}         % filler-page-format, e.g dots

\titlecontents{subsection}
    [1.0cm]             % left margin
    {\vspace{0.3cm}}                  % above code
    {\thecontentslabel \ } % numbered format
    {}         % unnumbered format
    {\titlerule*[0.3pc]{.}\contentspage}         % filler-page-format, e.g dots

 % Tạo mục lục tự động
\addtocontents{toc}{\protect\thispagestyle{empty}}
\tableofcontents
\thispagestyle{empty}
\cleardoublepage

\pagenumbering{roman}
%Tạo danh mục hình vẽ.
\renewcommand{\listfigurename}{DANH MỤC HÌNH VẼ}
{\let\oldnumberline\numberline
\renewcommand{\numberline}{Hình~\oldnumberline}
\listoffigures}
\newpage


 %Tạo danh mục bảng biểu.
\renewcommand{\listtablename}{DANH MỤC BẢNG BIỂU}
{\let\oldnumberline\numberline
\renewcommand{\numberline}{Bảng~\oldnumberline}
\listoftables}

\glsaddall
\renewcommand*{\acronymname}{DANH MỤC TỪ VIẾT TẮT}
\renewcommand*{\entryname}{Viết tắt}
\renewcommand*{\descriptionname}{Ý nghĩa}
\printnoidxglossaries

\newpage
\subfile{Chapter/0_5_Codes.tex}

\renewcommand\appendixname{PHỤ LỤC}
\renewcommand\appendixpagename{PHỤ LỤC}
\renewcommand\appendixtocname{PHỤ LỤC}

% ===================================================


\newpage
\pagenumbering{arabic}

\pagestyle{fancy}
\fancyhf{}
\fancyhead[RE, LO]{\small\leftmark} % \small: tên chương dài, cỡ thường bị xuống 2 dòng ở header
\fancyfoot[RE, LO]{\thepage}

\chapter{GIỚI THIỆU ĐỀ TÀI}
\label{chapter:Introduction}
\subfile{Chapter/1_Introduction}

\newpage
\chapter{NỀN TẢNG LÝ THUYẾT}
\label{chapter:Related_works}
\subfile{Chapter/2_Literature_review}

\newpage
% Phương pháp đề xuất chia thành hai chương (mẫu SoICT Research cho phép tách phương pháp
% thành nhiều chương): sinh luật và thực thi luật. Chương 3, 4 chỉ trình bày ý tưởng; nút giao cụ thể
% chỉ xuất hiện ở chương thực nghiệm. Chương PHÂN TÍCH LÝ THUYẾT của mẫu không bắt buộc và không dùng.
\chapter{SINH LUẬT HÀNH VI}
\label{chapter:RuleGeneration}
\subfile{Chapter/3_Rule_generation}

\newpage
\chapter{THỰC THI LUẬT HÀNH VI}
\label{chapter:RuleExecution}
\subfile{Chapter/4_Rule_execution}

\newpage
\chapter{THỰC NGHIỆM VÀ ĐÁNH GIÁ}
\label{chapter:Experiment}
\subfile{Chapter/5_Experiments}

\newpage
\chapter{KẾT LUẬN}
\label{chapter:conclusion}
\subfile{Chapter/6_Conclusions}

% ===================================================
\newpage
\renewcommand\bibname{TÀI LIỆU THAM KHẢO}
\printbibliography[heading=bibintoc] % bibintoc: mục lục ghi đúng trang đầu của danh mục

\newpage
\phantomsection\addappheadtotoc % Đặt trước \appendixpage để mục lục ghi đúng trang PHỤ LỤC
\appendixpage
\appendix

\titleformat{\chapter}[hang]{\centering\bfseries}{ \thechapter.\ }{0pt}{}[]
\titlespacing*{\chapter}{0pt}{-20pt}{20pt}

\titlecontents{chapter}
    [0.0cm]             % left margin
    {\bfseries\vspace{0.3cm}}                  % above code
    {{\bfseries{\scshape} \thecontentslabel.\ }} % numbered format
    {}         % unnumbered format
    {\titlerule*[0.3pc]{.}\contentspage}

\chapter{MIỀN THAM SỐ VÀ BỘ LUẬT B1}
\subfile{Chapter/Appendix_A}

\newpage
\chapter{LƯỢC ĐỒ DỮ LIỆU QUAN SÁT}
\subfile{Chapter/Appendix_B}

\newpage
\chapter{QUÁ TRÌNH XÂY DỰNG KỊCH BẢN}
\label{appendix:history}
\subfile{Chapter/Appendix_C}

\newpage
\chapter{PHẦN MỀM VÀ TÁI LẬP}
\label{appendix:software}
\subfile{Chapter/Appendix_D}

\end{document}
